% GENERATED FILE. Edit canonical lessons, then run npm run generate:books.
% canonical-source-hash: fnv1a64:a3eac6c619ddf342
% canonical-lessons: BN-C136-gobhir, BN-C136-choora, BN-C136-soru, BN-C136-gol, BN-C136-beguni

\chapter{Deep, Wide, Narrow, Round, Purple}
\label{ch:bn-deep-wide-narrow-round-purple}
\hlchaptermodality{136}

\begin{chapteropening}
\textbf{By the end of this chapter:} \emph{I can say deep, wide, narrow, round and purple.}

Complete the last lesson of chapter 136: I can say deep, wide, narrow, round and purple.
\end{chapteropening}

\section[\texorpdfstring{\bn{গভীর} (gôbhīr)}{গভীর (gôbhīr)}]{\bn{গভীর} (gôbhīr) — deep}

\label{lesson:BN-C136-gobhir}



\begin{quote}
Before the new one: say the Bengali for fast, then the Bengali for slow.
\end{quote}

\subsection*{You'll want to know: \bn{গভীর}}
\textbf{\bn{গভীর}} — \emph{gôbhīr} — \textquotedblleft{}deep\textquotedblright{}.

\begin{grammarlens}[title={What you've built}]
One more word for this chapter.
\end{grammarlens}

\subsection*{Your turn}
\begin{itemize}
\raggedright
  \item \emph{Say it:} \emph{gôbhīr}
  \item \emph{Say it:} \emph{gôbhīr}, once more
  \item \emph{Say it:} say \emph{gôbhīr}
  \item \emph{Recall:} say the Bengali for fast, then the Bengali for slow, then say \emph{gôbhīr} again
\end{itemize}

\begin{tcolorbox}[breakable,colback=teal!4,colframe=teal!35!black,title={Before you move on}]
What does \bn{গভীর} mean? (\textquotedblleft{}Deep\textquotedblright{}.) Say it once more.
\end{tcolorbox}
\section[\texorpdfstring{\bn{চওড়া} (chôoṛā)}{চওড়া (chôoṛā)}]{\bn{চওড়া} (chôoṛā) — wide}

\label{lesson:BN-C136-choora}



\begin{quote}
Before the new one: say the Bengali for slow, then the Bengali for deep.
\end{quote}

\subsection*{You'll want to know: \bn{চওড়া}}
\textbf{\bn{চওড়া}} — \emph{chôoṛā} — \textquotedblleft{}wide\textquotedblright{}.

\begin{grammarlens}[title={What you've built}]
One more word for this chapter.
\end{grammarlens}

\subsection*{Your turn}
\begin{itemize}
\raggedright
  \item \emph{Say it:} \emph{chôoṛā}
  \item \emph{Say it:} \emph{chôoṛā}, once more
  \item \emph{Say it:} say \emph{chôoṛā}
  \item \emph{Recall:} say the Bengali for slow, then the Bengali for deep, then say \emph{chôoṛā} again
\end{itemize}

\begin{tcolorbox}[breakable,colback=teal!4,colframe=teal!35!black,title={Before you move on}]
What does \bn{চওড়া} mean? (\textquotedblleft{}Wide\textquotedblright{}.) Say it once more.
\end{tcolorbox}
\section[\texorpdfstring{\bn{সরু} (sôru)}{সরু (sôru)}]{\bn{সরু} (sôru) — narrow}

\label{lesson:BN-C136-soru}



\begin{quote}
Before the new one: say the Bengali for deep, then the Bengali for wide.
\end{quote}

\subsection*{You'll want to know: \bn{সরু}}
\textbf{\bn{সরু}} — \emph{sôru} — \textquotedblleft{}narrow\textquotedblright{}.

\begin{grammarlens}[title={What you've built}]
One more word for this chapter.
\end{grammarlens}

\subsection*{Your turn}
\begin{itemize}
\raggedright
  \item \emph{Say it:} \emph{sôru}
  \item \emph{Say it:} \emph{sôru}, once more
  \item \emph{Say it:} say \emph{sôru}
  \item \emph{Recall:} say the Bengali for deep, then the Bengali for wide, then say \emph{sôru} again
\end{itemize}

\begin{tcolorbox}[breakable,colback=teal!4,colframe=teal!35!black,title={Before you move on}]
What does \bn{সরু} mean? (\textquotedblleft{}Narrow\textquotedblright{}.) Say it once more.
\end{tcolorbox}
\section[\texorpdfstring{\bn{গোল} (gol)}{গোল (gol)}]{\bn{গোল} (gol) — round}

\label{lesson:BN-C136-gol}



\begin{quote}
Before the new one: say the Bengali for wide, then the Bengali for narrow.
\end{quote}

\subsection*{You'll want to know: \bn{গোল}}
\textbf{\bn{গোল}} — \emph{gol} — \textquotedblleft{}round\textquotedblright{}.

\begin{grammarlens}[title={What you've built}]
One more word for this chapter.
\end{grammarlens}

\subsection*{Your turn}
\begin{itemize}
\raggedright
  \item \emph{Say it:} \emph{gol}
  \item \emph{Say it:} \emph{gol}, once more
  \item \emph{Say it:} say \emph{gol}
  \item \emph{Recall:} say the Bengali for wide, then the Bengali for narrow, then say \emph{gol} again
\end{itemize}

\begin{tcolorbox}[breakable,colback=teal!4,colframe=teal!35!black,title={Before you move on}]
What does \bn{গোল} mean? (\textquotedblleft{}Round\textquotedblright{}.) Say it once more.
\end{tcolorbox}
\section[\texorpdfstring{\bn{বেগুনি} (beguni)}{বেগুনি (beguni)}]{\bn{বেগুনি} (beguni) — purple}

\label{lesson:BN-C136-beguni}



\begin{quote}
Before the new one: say the Bengali for narrow, then the Bengali for round.
\end{quote}

\subsection*{You'll want to know: \bn{বেগুনি}}
\textbf{\bn{বেগুনি}} — \emph{beguni} — \textquotedblleft{}purple\textquotedblright{}.

\begin{grammarlens}[title={What you've built}]
One more word for this chapter.
\end{grammarlens}

\subsection*{Your turn}
\begin{itemize}
\raggedright
  \item \emph{Say it:} \emph{beguni}
  \item \emph{Say it:} \emph{beguni}, once more
  \item \emph{Say it:} say \emph{beguni}
  \item \emph{Recall:} say the Bengali for narrow, then the Bengali for round, then say \emph{beguni} again
\end{itemize}

\begin{tcolorbox}[breakable,colback=teal!4,colframe=teal!35!black,title={Before you move on}]
What does \bn{বেগুনি} mean? (\textquotedblleft{}Purple\textquotedblright{}.) Say it once more.
\end{tcolorbox}
